\begin{table}[H]
  \centering
  \small
  \caption{Méthode GIMSI adaptée dans un contexte Scrum.}
  \label{tab:ch02-gimsi-phases}
  \PFETableSetup
  % Tableau plus large que le bloc de texte : déborde de façon symétrique dans les marges.
  \makebox[\textwidth][c]{%
  \begingroup
  \renewcommand{\tabularxcolumn}[1]{>{\raggedright\arraybackslash}m{#1}}%
  \PFETableArrayStretch
  \arrayrulecolor{black}%
  \begin{tabularx}{\dimexpr\PFETableBlockWidth+2.8cm\relax}{%
    |>{\centering\arraybackslash}m{3.1cm}%
    |>{\centering\arraybackslash}m{0.55cm}%
    |>{\raggedright\arraybackslash}m{2.45cm}%
    |>{\raggedright\arraybackslash}X|}%
    \hline
    \tableHeaderStyle
    \tableHeaderCell{Phase} &
    \tableHeaderCell{N\textsuperscript{o}} &
    \tableHeaderCell{Étape} &
    \tableHeaderCell{Adaptation dans Scrum} \\
    \hline
    \multirow{2}{*}{\parbox{3.0cm}{\centering\bfseries Identification}} &
      \textbf{1} & Environnement de l’entreprise &
      Défini progressivement dans le Product Backlog. \\
    \cline{2-4}
    & \textbf{2} & Identification des acteurs &
      Réalisé via les User Stories et échanges avec les parties prenantes. \\
    \hline
    \multirow{5}{*}{\parbox{3.0cm}{\centering\bfseries Conception}} &
      \textbf{3} & Définition des objectifs &
      Affinés à chaque Sprint Planning. \\
    \cline{2-4}
    & \textbf{4} & Construction du tableau de bord &
      Conception itérative avec retour utilisateur. \\
    \cline{2-4}
    & \textbf{5} & Choix des indicateurs &
      KPI ajustés selon les besoins métiers évolutifs. \\
    \cline{2-4}
    & \textbf{6} & Collecte des informations &
      Données intégrées progressivement. \\
    \cline{2-4}
    & \textbf{7} & Système de tableaux de bord &
      Amélioration continue des tableaux de bord. \\
    \hline
    \multirow{2}{*}{\parbox{3.0cm}{\centering\bfseries Mise en œuvre}} &
      \textbf{8} & Choix des outils &
      Power BI retenu selon les besoins du projet. \\
    \cline{2-4}
    & \textbf{9} & Intégration &
      Réalisé par incréments à chaque sprint. \\
    \hline
    \parbox{3.0cm}{\centering\bfseries Suivi permanent} &
      \textbf{10} & Audit et amélioration &
      Amélioration continue via retours et itérations Scrum. \\
    \hline
  \end{tabularx}%
  \endgroup
  }%
\end{table}
